\section{Discussion}
\label{sec:discussion}

\paragraph{What II-A establishes.} Under a validated offline replay
model of a real 30B-A3B MoE coding workload, a substantial gap exists
between causal LRU-family retention and a non-causal Belady-optimal
upper bound, at every tested capacity-constrained expert-memory budget.
This gap -- the policy-information headroom -- is not uniform across
budgets: it is largest, among the tested points, near $8$~GiB in relative
terms and near $1.5$~GiB in absolute terms, and it collapses entirely
once the working set becomes effectively resident (the first tested such
budget being $16$~GiB). The gap persists after decomposing away the
specific cost of the runtime's sharded retention architecture, which is
small and not consistently signed across the sweep.

\paragraph{What II-B establishes.} Under the tested
Qwen3-30B-A3B-Q4\_K\_M protocol, routing is descriptively non-uniform,
sequentially dependent at one-step same-layer resolution, and temporally
local at short horizons relative to a shuffled null. At the same time,
R2-08 shows that cross-workload routing differences exceed temporal drift
under the predeclared comparison (ratio $2.21\times$), with
\texttt{WORKLOAD\_DEPENDENCE\_DOMINANT} as the regime classification.
These results characterize the routing stream; they do not evaluate a
policy.

\paragraph{Observed, inferred motivation, and not established.}
The distinction among these categories is central:
\begin{itemize}
 \item \textbf{Observed:} strong Oracle headroom (II-A); aggregate routing
   concentration, with LFU underperforming Global-LRU at the II-A anchor
   budget; measurable same-layer transition dependence (positive NMI in
   every tested layer); strong short-horizon reuse relative to a shuffled
   null; and cross-workload divergence exceeding temporal drift under the
   deterministic R2-08 protocol.
 \item \textbf{Inference/motivation:} these properties make strictly
   causal working-set prediction scientifically worth testing. The
   short-horizon locality is the most relevant signal for that test; the
   workload dependence implies that a practical predictor may need to be
   robust across workload categories or conditioned on workload context.
 \item \textbf{Not established:} that any causal predictor can recover a
   meaningful fraction of the Oracle-bounded policy-information headroom;
   that prediction accuracy is sufficient; that prefetch, admission, or a
   new retention policy improves runtime; that latency, throughput, or
   memory traffic improves; or that the observed structure generalizes
   across models, quantizations, machines, or non-deterministic decoding
   protocols.
\end{itemize}

\paragraph{Why II-A headroom is not automatically deployable.} An oracle
upper bound requires information no real system has at decision time.
Treating a large oracle gap as an implicit promise of recoverable
performance would conflate an informational upper bound with demonstrated
causal recovery. II-B adds structure, not recovery.

\paragraph{Why II-B structure is not automatically predictive.}
Concentration, transition dependence, reuse locality, and workload
dependence are descriptive properties of observed traces. Each can
motivate a hypothesis about future routing, but none is the same object as
a causal predictor evaluated on future decisions. This is why
Eq.~\eqref{eq:iib-iic-recovery} remains \notyetmeasured{} until
Research~II-C.